\chapter{The Deutsch--Jozsa and Bernstein--Vazirani Algorithms}

The algorithms in this chapter are historically the first to demonstrate, in a
clean and \emph{deterministic} way, that a quantum computer can outperform any
classical computer. The tasks they solve are artificial and of little practical
value, yet the tools they introduce---the oracle model, phase kickback, and
Hadamard interference---form the conceptual backbone of nearly every later
quantum algorithm, from Simon's problem to Shor's factoring algorithm. Master
these three ideas and the rest of the course rests on solid ground.

Throughout, $f$ denotes an unknown Boolean function that we may access only
through an oracle. Our measure of cost is the \emph{number of oracle queries},
not the number of elementary gates. We will see that a task requiring
exponentially many classical queries can be solved with a single quantum query.

\section{The Query and Oracle Model}

We are given a Boolean function $f:\{0,1\}^n\to\{0,1\}$ whose internal structure
is hidden from us. We do not have the code of $f$; we may only feed inputs and
receive outputs through a black box. To use $f$ inside a quantum circuit we need
its action to be reversible, since every quantum gate is unitary. The standard
device for this is the following.

\begin{definition}[Oracle / query model]
\label{l3:def:oracle}
Let $f:\{0,1\}^n\to\{0,1\}$. The \emph{oracle} (or \emph{quantum oracle}) for $f$
is the operator $\Uf$ acting on $n+1$ qubits---an $n$-qubit \emph{input register}
and a single-qubit \emph{target register}---defined on the computational basis by
\[
  \Uf\ket{x}\ket{y}=\ket{x}\ket{y\oplus f(x)},
  \qquad x\in\{0,1\}^n,\ y\in\{0,1\},
\]
where $\oplus$ denotes addition modulo $2$. The \emph{query complexity} of a task
is the minimum number of applications of $\Uf$ (over all algorithms that succeed)
needed to solve it.
\end{definition}

The exclusive-or in the target register---rather than simply overwriting $y$ with
$f(x)$---is precisely what makes the oracle reversible.

\begin{lemma}[The oracle is unitary]
\label{l3:lem:unitary}
For every $f:\{0,1\}^n\to\{0,1\}$, the operator $\Uf$ of Definition~\ref{l3:def:oracle}
is unitary and is its own inverse: $\Uf^{\dagger}\Uf=\Uf^2=\id$.
\end{lemma}

\begin{proof}
The set $\{\ket{x}\ket{y}\}$ over $x\in\{0,1\}^n$, $y\in\{0,1\}$ is an orthonormal
basis of the $2^{n+1}$-dimensional state space. The map
$(x,y)\mapsto(x,\,y\oplus f(x))$ is a bijection of the index set to itself, because
it is inverted by itself: applying it twice gives
$(x,\,y\oplus f(x)\oplus f(x))=(x,y)$ since $f(x)\oplus f(x)=0$. Hence $\Uf$ merely
permutes an orthonormal basis, so it is unitary, and $\Uf^2=\id$. Equivalently, in
matrix form $\braket{x'y'|\Uf^{\dagger}\Uf|xy}=\delta_{xx'}\delta_{yy'}$, i.e.
$\Uf^{\dagger}\Uf=\id$.
\end{proof}

\paragraph{Classical versus quantum access.}
A classical algorithm can only submit one input $x$ at a time and read back the
single bit $f(x)$; setting the target to $y=0$ gives $\Uf\ket{x}\ket{0}=\ket{x}\ket{f(x)}$.
To learn a \emph{global} property of $f$ this way may force us to examine almost
all $2^n$ inputs. A quantum algorithm can instead prepare a uniform superposition
of all inputs and apply $\Uf$ once, acting on every input in parallel:
\[
  \Uf\Big(\tfrac{1}{\sqrt{N}}\sum_{x\in\{0,1\}^n}\ket{x}\Big)\ket{0}
  =\tfrac{1}{\sqrt{N}}\sum_{x}\ket{x}\ket{f(x)},
  \qquad N=2^n.
\]
Parallelism alone is not enough---measuring this state returns one random pair
$(x,f(x))$, no better than a classical query. The extra ingredient that converts
parallelism into advantage is \emph{interference}, engineered through phase
kickback and the Hadamard transform.

\section{Phase Kickback}

The oracle writes $f(x)$ into the target register as a bit. We would prefer to
have $f(x)$ appear as a \emph{phase} on the input register, since phases are what
interfere. The following observation---perhaps the single most important trick in
the chapter---accomplishes exactly this by choosing the target register cleverly.

\begin{lemma}[Phase kickback]
\label{l3:lem:kickback}
Let $\ket{-}=\tfrac{1}{\sqrt2}(\ket{0}-\ket{1})$. Then for every $x\in\{0,1\}^n$,
\[
  \Uf\ket{x}\ket{-}=(-1)^{f(x)}\ket{x}\ket{-}.
\]
In particular the target qubit is left in $\ket{-}$ and may be discarded, so on
the input register the oracle acts as the effective diagonal unitary
$\tilde U_f\ket{x}=(-1)^{f(x)}\ket{x}$.
\end{lemma}

\begin{proof}
Expand the target register and apply Definition~\ref{l3:def:oracle} to each branch:
\[
  \Uf\ket{x}\ket{-}
  =\tfrac{1}{\sqrt2}\big(\Uf\ket{x}\ket{0}-\Uf\ket{x}\ket{1}\big)
  =\tfrac{1}{\sqrt2}\big(\ket{x}\ket{f(x)}-\ket{x}\ket{1\oplus f(x)}\big).
\]
Consider the two cases for the bit $f(x)$. If $f(x)=0$ the target factor is
$\tfrac{1}{\sqrt2}(\ket{0}-\ket{1})=\ket{-}=(+1)\ket{-}$. If $f(x)=1$ it is
$\tfrac{1}{\sqrt2}(\ket{1}-\ket{0})=-\ket{-}=(-1)\ket{-}$. Both cases are captured
by the single formula $(-1)^{f(x)}\ket{-}$, whence
$\Uf\ket{x}\ket{-}=(-1)^{f(x)}\ket{x}\ket{-}$. Since the target factor is unchanged
and unentangled from the input register, it factors out and can be dropped.
\end{proof}

\begin{remark}
The name \emph{kickback} reflects that the eigenvalue $(-1)^{f(x)}$ of the target
qubit's controlled action is ``kicked back'' onto the control (input) register.
The target qubit acts purely as a passive signal state. This is the same mechanism
that, with a target eigenstate of a general unitary, powers quantum phase
estimation later in the course.
\end{remark}

Applying $\Uf$ to a uniform superposition with the target in $\ket{-}$ therefore
produces, on the input register,
\[
  \tfrac{1}{\sqrt{N}}\sum_{x}\ket{x}\;\xmapsto{\ \Uf\ }\;
  \tfrac{1}{\sqrt{N}}\sum_{x}(-1)^{f(x)}\ket{x}.
\]
Every value of $f$ is now recorded in a relative sign. It remains to read out
these signs; the tool for that is the Hadamard transform.

\section{The Hadamard Transform}

We will use repeatedly the action of $H^{\otimes n}$ on a basis state. Recall for a
single qubit $H\ket{a}=\tfrac{1}{\sqrt2}\sum_{b\in\{0,1\}}(-1)^{ab}\ket{b}$.

\begin{lemma}[Hadamard identity]
\label{l3:lem:had}
For $x\in\{0,1\}^n$, writing $x\cdot z=\bigoplus_{i=1}^{n}x_i z_i$ for the bitwise
inner product modulo $2$,
\[
  H^{\otimes n}\ket{x}=\frac{1}{\sqrt{2^{\,n}}}\sum_{z\in\{0,1\}^n}(-1)^{x\cdot z}\ket{z}.
\]
\end{lemma}

\begin{proof}
Take the tensor product of the single-qubit rule over the $n$ qubits:
\[
  H^{\otimes n}\ket{x}
  =\bigotimes_{i=1}^{n}\Big(\tfrac{1}{\sqrt2}\sum_{z_i\in\{0,1\}}(-1)^{x_i z_i}\ket{z_i}\Big)
  =\frac{1}{\sqrt{2^{\,n}}}\sum_{z\in\{0,1\}^n}(-1)^{\sum_i x_i z_i}\ket{z}.
\]
Because the phase is $\pm1$, the exponent $\sum_i x_i z_i$ may be reduced modulo
$2$, giving $(-1)^{x\cdot z}$.
\end{proof}

We also record the orthogonality relation that makes the interference exact:
\[
  \frac{1}{2^{\,n}}\sum_{x\in\{0,1\}^n}(-1)^{x\cdot w}
  =\delta_{w,0^n},
  \qquad w\in\{0,1\}^n,
\]
which holds because the sum factors as $\prod_{i=1}^{n}\tfrac12\sum_{x_i}(-1)^{x_i w_i}$,
and each factor equals $1$ if $w_i=0$ and $0$ if $w_i=1$.

\section{The Deutsch Algorithm}

Consider the simplest nontrivial case, $n=1$, so $f:\{0,1\}\to\{0,1\}$. There are
four such functions. We wish to decide a global property of $f$ using as few
queries as possible.

\begin{definition}[Constant and balanced functions]
\label{l3:def:constbal}
A function $f:\{0,1\}^n\to\{0,1\}$ is \emph{constant} if $f(x)$ takes the same value
on all $2^n$ inputs. It is \emph{balanced} if $f(x)=0$ on exactly half of the inputs
and $f(x)=1$ on the other half. For $n=1$, $f$ is constant iff $f(0)=f(1)$ and
balanced iff $f(0)\neq f(1)$.
\end{definition}

Classically, deciding whether $f$ is constant or balanced for $n=1$ requires
\emph{two} queries: knowing $f(0)$ alone tells us nothing about $f(1)$. Deutsch's
algorithm decides it with a single query.

\begin{theorem}[Correctness of the Deutsch algorithm]
\label{l3:thm:deutsch}
Prepare the two-qubit state $\ket{0}\ket{1}$, apply $H$ to each qubit, apply $\Uf$,
apply $H$ to the first qubit, and measure the first qubit. The outcome is $0$ if
$f$ is constant and $1$ if $f$ is balanced. Exactly one query to $\Uf$ is used.
\end{theorem}

\begin{proof}
After the first Hadamards,
$H\ket{0}\otimes H\ket{1}=\ket{+}\ket{-}=\tfrac{1}{\sqrt2}(\ket{0}+\ket{1})\ket{-}$.
Applying $\Uf$ and using phase kickback (Lemma~\ref{l3:lem:kickback}) on each branch,
\[
  \Uf\,\tfrac{1}{\sqrt2}\big(\ket{0}+\ket{1}\big)\ket{-}
  =\tfrac{1}{\sqrt2}\big((-1)^{f(0)}\ket{0}+(-1)^{f(1)}\ket{1}\big)\ket{-}.
\]
Now apply $H$ to the first qubit, using $H\ket{0}=\tfrac{1}{\sqrt2}(\ket0+\ket1)$
and $H\ket{1}=\tfrac{1}{\sqrt2}(\ket0-\ket1)$:
\[
  \tfrac{1}{2}\Big[\big((-1)^{f(0)}+(-1)^{f(1)}\big)\ket{0}
  +\big((-1)^{f(0)}-(-1)^{f(1)}\big)\ket{1}\Big]\ket{-}.
\]
If $f$ is constant, $f(0)=f(1)$, the $\ket{1}$ amplitude vanishes and the
$\ket{0}$ amplitude is $\pm1$; the first qubit is $\pm\ket{0}$ and measurement
yields $0$ with certainty. If $f$ is balanced, $f(0)\neq f(1)$, the $\ket{0}$
amplitude vanishes and the first qubit is $\pm\ket{1}$; measurement yields $1$
with certainty. The oracle is called once.
\end{proof}

Notice that the answer is read from the \emph{first} qubit, even though classically
one would expect the ``information'' to live in the target qubit that received
$f(x)$. This inversion is the hallmark of interference.

\begin{example}[Deutsch on $f(a)=a\oplus 1$]
\label{l3:ex:deutsch}
Take $f(0)=1$, $f(1)=0$, which is balanced. After the query the first register is
$\tfrac{1}{\sqrt2}\big((-1)^{1}\ket0+(-1)^{0}\ket1\big)=\tfrac{1}{\sqrt2}(-\ket0+\ket1)=-\ket{-}$.
Applying $H$ gives $H(-\ket{-})=-\ket{1}$, so the measurement returns $1$,
correctly reporting ``balanced.'' By contrast, for the constant function
$f\equiv 1$ the first register after the query is
$\tfrac{1}{\sqrt2}(-\ket0-\ket1)=-\ket{+}$, and $H(-\ket{+})=-\ket0$ returns $0$,
reporting ``constant.''
\end{example}

\section{The Deutsch--Jozsa Algorithm}

We now generalize to $n$ qubits under the following \emph{promise}: $f$ is either
constant or balanced (Definition~\ref{l3:def:constbal}), and we are guaranteed that
one of these holds. The task is to decide which. Classically, in the worst case
one must query $2^{n-1}+1$ inputs: after seeing $2^{n-1}$ identical outputs the
function could still be either constant or balanced, and one more query is needed.
Deutsch--Jozsa decides it with a single quantum query.

The circuit is the unifying pattern $H^{\otimes n}$--$\Uf$--$H^{\otimes n}$ on the
$n$-qubit input register, with the target qubit held in $\ket{-}$.

\begin{theorem}[Correctness of Deutsch--Jozsa]
\label{l3:thm:dj}
Prepare $\ket{0}^{\otimes n}$ on the input register and $\ket{-}$ on the target.
Apply $H^{\otimes n}$, then $\Uf$, then $H^{\otimes n}$, and measure the input
register. If $f$ is constant the outcome is $0^n$ with probability $1$; if $f$ is
balanced the outcome $0^n$ occurs with probability $0$. Thus measuring all zeros
certifies ``constant,'' and any nonzero outcome certifies ``balanced.'' One query
to $\Uf$ is used.
\end{theorem}

\begin{proof}
By Lemma~\ref{l3:lem:had}, $H^{\otimes n}\ket{0}^{\otimes n}
=\tfrac{1}{\sqrt N}\sum_{x}\ket{x}$ with $N=2^n$. Applying $\Uf$ with the target
in $\ket{-}$ triggers phase kickback (Lemma~\ref{l3:lem:kickback}) on every branch,
giving (dropping the unchanged target factor)
\[
  \tfrac{1}{\sqrt N}\sum_{x}(-1)^{f(x)}\ket{x}.
\]
Apply $H^{\otimes n}$ once more and use Lemma~\ref{l3:lem:had} again:
\[
  H^{\otimes n}\tfrac{1}{\sqrt N}\sum_{x}(-1)^{f(x)}\ket{x}
  =\frac{1}{N}\sum_{z}\Big(\sum_{x}(-1)^{f(x)}(-1)^{x\cdot z}\Big)\ket{z}.
\]
The amplitude of the outcome $z=0^n$ is obtained by setting $z=0^n$, so $x\cdot z=0$:
\[
  \alpha_{0^n}=\frac{1}{N}\sum_{x\in\{0,1\}^n}(-1)^{f(x)}.
\]
If $f$ is constant, every term is the same sign and $\alpha_{0^n}=\pm1$, so the
outcome is $0^n$ with probability $\abs{\alpha_{0^n}}^2=1$. If $f$ is balanced,
exactly $N/2$ terms equal $+1$ and $N/2$ equal $-1$, so they cancel and
$\alpha_{0^n}=0$: the outcome $0^n$ never occurs. In either case a single query
suffices, and the two cases are distinguished with certainty.
\end{proof}

\section{The Bernstein--Vazirani Algorithm}

The Deutsch--Jozsa algorithm returns a single bit of information (constant vs.\
balanced). The Bernstein--Vazirani algorithm uses the very same circuit to extract
an entire hidden $n$-bit string in one query, provided $f$ has a special linear
structure.

\begin{definition}[Bernstein--Vazirani function]
\label{l3:def:bv}
Fix a hidden string $s\in\{0,1\}^n$. The associated \emph{linear} Boolean function is
\[
  f_s(x)=s\cdot x=\bigoplus_{i=1}^{n}s_i x_i \pmod 2,
\]
the parity of the number of positions at which $s$ and $x$ are both $1$. The
Bernstein--Vazirani problem is: given oracle access to $\Uf$ for such an $f_s$,
determine $s$.
\end{definition}

Classically, one recovers $s$ bit by bit. Querying $x=e_i$ (the string with a
single $1$ in position $i$) returns $f_s(e_i)=s_i$; this determines one bit per
query, so $n$ queries are needed and sufficient. The quantum algorithm needs one.

\begin{theorem}[Correctness of Bernstein--Vazirani]
\label{l3:thm:bv}
For $f=f_s$ of Definition~\ref{l3:def:bv}, running the circuit of
Theorem~\ref{l3:thm:dj}---$H^{\otimes n}$, then $\Uf$, then $H^{\otimes n}$, with the
target in $\ket{-}$---leaves the input register in the state $\ket{s}$. Measuring
recovers $s$ with certainty using a single query.
\end{theorem}

\begin{proof}
As before, after $H^{\otimes n}$ and the single query with phase kickback the input
register is
\[
  \tfrac{1}{\sqrt N}\sum_{x}(-1)^{f_s(x)}\ket{x}
  =\tfrac{1}{\sqrt N}\sum_{x}(-1)^{s\cdot x}\ket{x}.
\]
Applying $H^{\otimes n}$ and using Lemma~\ref{l3:lem:had},
\[
  H^{\otimes n}\tfrac{1}{\sqrt N}\sum_{x}(-1)^{s\cdot x}\ket{x}
  =\frac{1}{N}\sum_{z}\Big(\sum_{x}(-1)^{s\cdot x}(-1)^{x\cdot z}\Big)\ket{z}
  =\frac{1}{N}\sum_{z}\Big(\sum_{x}(-1)^{x\cdot(s\oplus z)}\Big)\ket{z},
\]
where $s\cdot x\oplus x\cdot z=x\cdot(s\oplus z)$ because the inner product is
linear in its second argument modulo $2$. By the orthogonality relation of
Section~3, $\tfrac{1}{N}\sum_{x}(-1)^{x\cdot(s\oplus z)}=\delta_{s\oplus z,\,0^n}
=\delta_{z,s}$. Hence only the term $z=s$ survives and the register is exactly
$\ket{s}$. Measuring in the computational basis returns $s$ with probability $1$.
\end{proof}

Equivalently, $H^{\otimes n}\ket{s}=\tfrac{1}{\sqrt N}\sum_x(-1)^{s\cdot x}\ket{x}$
by Lemma~\ref{l3:lem:had}, so the post-query state is $H^{\otimes n}\ket{s}$; a second
$H^{\otimes n}$ inverts it (since $H^{\otimes n}$ is its own inverse) and returns
$\ket{s}$.

\begin{example}[A small Bernstein--Vazirani instance]
\label{l3:ex:bv}
Let $n=3$ and $s=101$, so $f_s(x)=x_1\oplus x_3$. After $H^{\otimes 3}$ and one
query, the input register is $\tfrac{1}{\sqrt8}\sum_{x}(-1)^{x_1\oplus x_3}\ket{x}$.
Applying $H^{\otimes 3}$ collapses this to $\ket{101}$, and the measurement reads
off $s=101$ in one shot. A classical solver would instead query
$x=100$, $010$, $001$ to obtain $s_1=1$, $s_2=0$, $s_3=1$ respectively---three
queries for the same answer.
\end{example}

\section{Discussion and Significance}

\begin{remark}[Deterministic advantage and its limits]
\label{l3:rem:significance}
Several points deserve emphasis.
\begin{itemize}
\item \emph{Deterministic vs.\ bounded-error.} All three algorithms succeed with
probability exactly $1$; the advantage is not a matter of a lucky measurement. For
Deutsch--Jozsa the separation is $1$ quantum query versus $2^{n-1}+1$ classical
queries in the worst case (deterministically). If one allows a small classical
error probability, a handful of random queries distinguishes constant from
balanced with high confidence, so the \emph{bounded-error} classical query
complexity of Deutsch--Jozsa is only $O(1)$; the exponential gap there is specific
to the deterministic setting. For Bernstein--Vazirani the separation is $1$ versus
$n$ queries and survives even in the bounded-error model.
\item \emph{The unifying pattern.} Every algorithm here is the circuit
$H^{\otimes n}\,\Uf\,H^{\otimes n}$ applied to $\ket{0}^{\otimes n}$ with the target
in $\ket{-}$. Phase kickback writes $f$ into signs; the two Hadamard layers read
those signs out by interference. This template recurs throughout quantum
algorithms.
\item \emph{Query advantage $\neq$ gate-complexity advantage.} We count calls to
$\Uf$, treating the oracle as a unit-cost black box. The number of elementary
gates needed to \emph{build} $\Uf$ for a given $f$ must be analyzed separately, and
for a generic $f$ it can be large. The dramatic query separations above do not by
themselves imply a separation in total running time.
\item \emph{Relation to Simon and Shor.} Bernstein--Vazirani generalizes to
Simon's problem, where $f$ obeys the hidden-structure promise
$f(x)=f(x\oplus s)$ (equivalently $f(x)=f(y)$ iff $x\oplus y\in\{0^n,s\}$); the
same $H^{\otimes n}\,\Uf\,H^{\otimes n}$ pattern yields linear equations
$z\cdot s=0$ constraining $s$, and Simon's exponential separation holds even with
bounded error. Simon's algorithm in turn inspired Shor's factoring algorithm,
where the final Hadamard transform is replaced by the quantum Fourier transform of
a later chapter.
\end{itemize}
\end{remark}

\section{Problems}

\begin{exercise}
\label{l3:ex:p1}
List all four functions $f:\{0,1\}\to\{0,1\}$ and classify each as constant or
balanced. For each, compute by hand the first-qubit measurement outcome produced
by the Deutsch algorithm of Theorem~\ref{l3:thm:deutsch}.
\end{exercise}

\begin{exercise}
\label{l3:ex:p2}
Prove that $\Uf$ of Definition~\ref{l3:def:oracle} is Hermitian, and use this together
with $\Uf^2=\id$ to conclude that its eigenvalues are $\pm1$.
\end{exercise}

\begin{exercise}
\label{l3:ex:p3}
Show directly, without invoking Lemma~\ref{l3:lem:kickback}, that if the target qubit
is initialized in $\ket{+}=\tfrac{1}{\sqrt2}(\ket0+\ket1)$ instead of $\ket{-}$,
then $\Uf\ket{x}\ket{+}=\ket{x}\ket{+}$ for all $x$, so no phase is kicked back.
Explain why $\ket{-}$ is the useful choice.
\end{exercise}

\begin{exercise}
\label{l3:ex:p4}
For $n=2$, let $f(x_1x_2)=x_1\oplus x_2$. Verify that $f$ is balanced and trace the
Deutsch--Jozsa state through $H^{\otimes 2}$, $\Uf$, $H^{\otimes 2}$, confirming
that the amplitude of $\ket{00}$ is $0$.
\end{exercise}

\begin{exercise}
\label{l3:ex:p5}
In the Bernstein--Vazirani setting with hidden string $s$, suppose the oracle is
for $g(x)=s\cdot x\oplus c$ with an unknown constant bit $c$ (an ``affine''
function). Show that the algorithm of Theorem~\ref{l3:thm:bv} still outputs $s$ with
certainty, and identify where the constant $c$ goes.
\end{exercise}

\begin{exercise}
\label{l3:ex:p6}
Prove the orthogonality relation $\tfrac{1}{2^n}\sum_{x\in\{0,1\}^n}(-1)^{x\cdot w}
=\delta_{w,0^n}$ used in the proof of Theorem~\ref{l3:thm:bv}, by factoring the sum
over the $n$ bits.
\end{exercise}

\begin{exercise}
\label{l3:ex:p7}
Give the exact worst-case classical (deterministic) query complexity of the
Deutsch--Jozsa problem on $n$ bits and justify it. Contrast this with the
observation that a randomized classical algorithm distinguishes constant from
balanced with error at most $1/3$ using only a constant number of queries.
\end{exercise}

\section{Solutions}

\begin{solution}[Exercise~\ref{l3:ex:p1}]
The four functions are: $f\equiv0$ (constant), $f\equiv1$ (constant),
$f(a)=a$ (balanced, $f(0)=0,f(1)=1$), and $f(a)=a\oplus1$ (balanced,
$f(0)=1,f(1)=0$). By Theorem~\ref{l3:thm:deutsch} the outcome is $0$ for the two
constant functions and $1$ for the two balanced ones. Concretely, for $f\equiv0$
the post-query first register is $\ket{+}\to H\ket{+}=\ket0$; for $f\equiv1$ it is
$-\ket{+}\to-\ket0$ (outcome $0$); for $f(a)=a$ it is
$\tfrac{1}{\sqrt2}(\ket0-\ket1)=\ket{-}\to\ket1$; for $f(a)=a\oplus1$ it is
$-\ket{-}\to-\ket1$ (outcome $1$).
\end{solution}

\begin{solution}[Exercise~\ref{l3:ex:p2}]
In the computational basis $\Uf$ is a real permutation matrix
($\ket{x}\ket{y}\mapsto\ket{x}\ket{y\oplus f(x)}$). Because the permutation is an
involution (Lemma~\ref{l3:lem:unitary}), the matrix equals its own transpose, hence is
symmetric; being real and symmetric it is Hermitian, $\Uf^{\dagger}=\Uf$.
Alternatively, from $\Uf^{\dagger}\Uf=\id$ (unitary) and $\Uf^2=\id$ (involution) we
get $\Uf^{\dagger}=\Uf^{-1}=\Uf$. If $\Uf\ket{\psi}=\lambda\ket{\psi}$ then
$\ket{\psi}=\Uf^2\ket{\psi}=\lambda^2\ket{\psi}$, so $\lambda^2=1$ and
$\lambda\in\{+1,-1\}$.
\end{solution}

\begin{solution}[Exercise~\ref{l3:ex:p3}]
Expanding as in Lemma~\ref{l3:lem:kickback},
\[
  \Uf\ket{x}\ket{+}
  =\tfrac{1}{\sqrt2}\big(\ket{x}\ket{f(x)}+\ket{x}\ket{1\oplus f(x)}\big)
  =\tfrac{1}{\sqrt2}\ket{x}\big(\ket{f(x)}+\ket{1\oplus f(x)}\big).
\]
Whether $f(x)=0$ or $f(x)=1$, the set $\{f(x),\,1\oplus f(x)\}$ equals $\{0,1\}$,
so the target factor is always $\tfrac{1}{\sqrt2}(\ket0+\ket1)=\ket{+}$ with a
$+$ sign. Thus $\Uf\ket{x}\ket{+}=\ket{x}\ket{+}$ and no phase depends on $f(x)$:
$\ket{+}$ is a $+1$ eigenstate of the target's bit-flip action. The state $\ket{-}$
is the $-1$ eigenstate, so flipping (which happens exactly when $f(x)=1$) multiplies
it by $-1$; that sign is what encodes $f$. Hence $\ket{-}$ is the useful choice.
\end{solution}

\begin{solution}[Exercise~\ref{l3:ex:p4}]
Here $f(00)=0,\ f(01)=1,\ f(10)=1,\ f(11)=0$: two zeros and two ones, so $f$ is
balanced. Starting from $\ket{00}$, $H^{\otimes2}$ gives
$\tfrac12(\ket{00}+\ket{01}+\ket{10}+\ket{11})$. Phase kickback multiplies each
$\ket{x}$ by $(-1)^{f(x)}$:
$\tfrac12(\ket{00}-\ket{01}-\ket{10}+\ket{11})$. Note this factors as
$\tfrac12(\ket0-\ket1)(\ket0-\ket1)=\ket{-}\ket{-}$. Applying $H^{\otimes2}$ gives
$H\ket{-}\otimes H\ket{-}=\ket1\ket1=\ket{11}$. The amplitude of $\ket{00}$ is $0$,
consistent with Theorem~\ref{l3:thm:dj}; indeed
$\alpha_{00}=\tfrac14\sum_x(-1)^{f(x)}=\tfrac14(1-1-1+1)=0$.
\end{solution}

\begin{solution}[Exercise~\ref{l3:ex:p5}]
With $g(x)=s\cdot x\oplus c$, phase kickback gives each branch the sign
$(-1)^{g(x)}=(-1)^{c}(-1)^{s\cdot x}$. The factor $(-1)^c$ is a \emph{global} phase
independent of $x$, so it can be pulled out of the sum:
\[
  \tfrac{1}{\sqrt N}\sum_x(-1)^{g(x)}\ket{x}
  =(-1)^{c}\,\tfrac{1}{\sqrt N}\sum_x(-1)^{s\cdot x}\ket{x}.
\]
Applying $H^{\otimes n}$ yields $(-1)^{c}\ket{s}$ by Theorem~\ref{l3:thm:bv}. A global
phase is physically undetectable and does not affect measurement probabilities, so
the outcome is $s$ with certainty. The constant $c$ is invisible to this algorithm;
recovering $c$ would require an additional query, e.g.\ evaluating $g(0^n)=c$.
\end{solution}

\begin{solution}[Exercise~\ref{l3:ex:p6}]
The inner product splits over bits, so the sum factors:
\[
  \frac{1}{2^n}\sum_{x\in\{0,1\}^n}(-1)^{x\cdot w}
  =\prod_{i=1}^{n}\Big(\tfrac12\sum_{x_i\in\{0,1\}}(-1)^{x_i w_i}\Big).
\]
For each factor, if $w_i=0$ the summand is $(-1)^0=1$ twice, giving
$\tfrac12(1+1)=1$; if $w_i=1$ the summand is $+1$ and $-1$, giving
$\tfrac12(1-1)=0$. The product is $1$ only when every $w_i=0$, i.e.\ $w=0^n$, and
$0$ otherwise. Hence the sum equals $\delta_{w,0^n}$.
\end{solution}

\begin{solution}[Exercise~\ref{l3:ex:p7}]
Deterministically, an adversary answering queries can keep all responses equal
(say all $0$) for the first $2^{n-1}$ distinct inputs; at that point $f$ could still
be constant (all outputs $0$) or balanced (the remaining $2^{n-1}$ inputs all $1$).
Only a $(2^{n-1}+1)$-th query can force the decision, and $2^{n-1}+1$ queries always
suffice: if any two of them disagree, $f$ is balanced, otherwise the majority
argument forces it to be constant. So the exact worst-case deterministic query
complexity is $2^{n-1}+1$. Randomly, sampling $k$ distinct inputs and answering
``balanced'' iff any two disagree: a constant $f$ is never misclassified, while a
balanced $f$ is missed only if all $k$ samples share a value, an event of
probability $2\cdot 2^{-k}$ (up to the small correction for sampling without
replacement). Taking $k=3$ already gives error below $1/3$, so a constant number of
queries suffices with bounded error---this is why the exponential separation is
special to the deterministic model, as noted in Remark~\ref{l3:rem:significance}.
\end{solution}
